\input{../tex-inputs/latex-leseansicht-vorspann}

\section[Arthur Schnitzler an Gustav Schwarzkopf, 7. 7. 1916]{L04403 Arthur Schnitzler an Gustav Schwarzkopf, 7. 7. 1916}
\nopagebreak\mylabel{L04403v}
\rehead{ }\normalsize\beginnumbering\briefempfaengerindex{Schwarzkopf, Gustav@\textsc{Schwarzkopf, Gustav}!zzzSchnitzler, Arthur@\emph{von Arthur Schnitzler}!1916-07-071@{7. 7. 1916}|(be}
\toendnotes[C]{\smallbreak\pagebreak[2]}
\correspDesc{Versand  durch Arthur Schnitzler am 7. 7. 1916 in Altaussee
\newline{}Übermittlung  am 7. 7. 1916 in Altaussee
\newline{}Erhalt  durch Gustav Schwarzkopf im Zeitraum [8. 7. 1916 – 12. 7. 1916?] in Marienbad}\toendnotes[C]{\smallbreak}
\Standort{DLA, A:Schnitzler, HS.1985.1.1897.}
\physDesc{Bildpostkarte, 556 Zeichen
\newline{}Handschrift: Bleistift, deutsche Kurrent
\newline{}Versand: Stempel: »\nobreak{}\oindex{Altaussee@\textbf{Altaussee}|pwk}Alt-Aussee, 7. VII. 1\textcolor{gray}{6}\nobreak{}«.  }\toendnotes[C]{\smallbreak}\pstart{}{\pb}Herrn Gustav Schwarzkopf\pend{}\pstart{}Wien I\oindex{I., Innere Stadt@\textbf{I., Innere Stadt}|pw}\pend{}\pstart{}Tiefer Graben 17\oindex{Wien@\textbf{Wien}!I., Innere Stadt@\textbf{I., Innere Stadt}!Tiefer Graben 17@\textbf{Tiefer Graben 17}|pw}\pend{}{\bigskip}
\pstart
           \centering{}\textcolor{gray}{\textbf{{\pb}Salzkammergut\oindex{Salzkammergut@\textbf{Salzkammergut}|pw}.}}\pend
           
\pstart
           \centering{}\textcolor{gray}{\textbf{Alt-Aussee\oindex{Altaussee@\textbf{Altaussee}|pw} mit dem Loser\oindex{Loser@\textbf{Loser}|pw}.}}\pend
           \vspace{1em}
\pstart
           \raggedleft{}{\pb}7. 7. 16\pend
           \vspace{0.5em}
\pstart
           lieber Guſtav, wir haben uns in Haus\oindex{Villa Annerl@\textbf{Villa Annerl}|pwv}
               und Landſchaft beſtens inſtallirt; – Verpflanzung vorläufig ohne jede Schwierigkeit,
               war auch nicht billig; (übrigens eſſen wir Mittags{ }ſehr gut beim Seewirth\oindex{Hotel am See@\textbf{Hotel am See}|pw} – und haben für die andern Mahlzeiten das meiſte ausWien\oindex{Wien@\textbf{Wien}|pw} mit).\pend
           
\pstart
           – \label{K_L04403-1v}\edtext{Geſtern}{\lemma{\textnormal{\emph{Gestern}}}\Cendnote{\textnormal{Vgl. A. S.: \emph{Tagebuch}, 6. 7. 1916.}}}\label{K_L04403-1} gab es gab es Geſang und Clavier für uns;
               Frau Hanſi Landesberger\pwindex{Nemetschke, Johanna 11.\,8.\,1882 Wien – 8.\,12.\,1958@\textsc{Nemetschke, Johanna} (11.\,8.\,1882 Wien – 8.\,12.\,1958)|pw}, ihre Tochter\pwindex{Landesberger, Susanne von 8.\,8.\,1910 – 14.\,12.\,2003 Rom@\textsc{Landesberger, Susanne von} (8.\,8.\,1910 – 14.\,12.\,2003 Rom)|pwv};  Frl v Filtſch\pwindex{Filtsch, Molly von 1855 Sibiu – 21.\,3.\,1926 Wien@\textsc{Filtsch, Molly von} (1855 Sibiu – 21.\,3.\,1926 Wien)|pw} (die ja auch eine
               Ihrer intimen Freundi{\geminationn}en zu{ }ſein{ }ſcheint, \textcolor{gray}{S}ie
               \textcolor{gray}{Weltma{\geminationn}}) {\pb}hörte zu. – Es
               wäre ſehr ſchön, we{\geminationn}{ }ſie herkommen.\pend
           
\pstart
           Wir grüßen herzlichſst{\\[\baselineskip]}Ihr \spacefill\mbox{Arthur}\pend
           \leftskip=0em{}\selectlanguage{ngerman}\endnumbering\briefempfaengerindex{Schwarzkopf, Gustav@\textsc{Schwarzkopf, Gustav}!zzzSchnitzler, Arthur@\emph{von Arthur Schnitzler}!1916-07-071@{7. 7. 1916}|)be}\mylabel{L04403h}
\begin{anhang}
\end{anhang}\newcommand{\dateiname}{L04403}\newcommand{\titel}{Arthur Schnitzler an Gustav Schwarzkopf, 7. 7. 1916}\newcommand{\editorInnen}{Herausgegeben von Jahnke, SelmaMüller, Martin Anton}\input{../tex-inputs/latex-leseansicht-abspann}
